\begin{thebibliography}{99}
\setlength{\itemsep}{2pt}
\bibitem{ALW}D. Aasen, E. Lake and K. Walker, \emph{Fermion condensation and super pivotal categories}, arXiv:1709.01941v2 (2018).

\bibitem{Integral}M. A. Alekseyev, W. Bruns, J. Dong and S. Palcoux, \emph{Classifying integral Grothendieck rings up to rank 5 and beyond}, arXiv:2507.07023 (2025).

\bibitem{AGM}I. Angiono, C. Galindo and G. Mora, \emph{Zestings of Hopf algebras}, arXiv:2505.10447v1 (2025).

\bibitem{BMPS}S. Bigelow, S. Morrison, E. Peters and N. Snyder,
\emph{Constructing the extended Haagerup planar algebra}, Acta Math.
\textbf{209} (2012), 29--82; arXiv:0909.4099.

\bibitem{BNRW}P. Bruillard, S.-H. Ng, E. C. Rowell and Z. Wang, \emph{Rank-finiteness for modular categories}, J. Amer. Math. Soc. \textbf{29} (2016), 857--881; arXiv:1310.7050v3.

\bibitem{BPEH}S. Burciu and S. Palcoux,
\emph{Structure constants, Isaacs property and extended Haagerup fusion categories},
Comm. Algebra \textbf{53} (2025), 1438--1452; arXiv:2210.13936v3,
especially Section~5.1.

\bibitem{CLS}Y. Choi, D.-C. Lu and Z. Sun, \emph{Self-duality under gauging a non-invertible symmetry}, J. High Energy Phys. \textbf{2024}, no.~1, article 142; arXiv:2310.19867v2.

\bibitem{CRT}R. Coquereaux, R. Rais and E. H. Tahri, \emph{Exceptional quantum subgroups for the rank two Lie algebras $B_2$ and $G_2$}, J. Math. Phys. \textbf{51} (2010), 092302; arXiv:1001.5416v1.

\bibitem{DGP}C. Delaney, C. Galindo, J. Y. Plavnik, E. C. Rowell and Q. Zhang, \emph{Braided zesting and its applications}, Comm. Math. Phys. \textbf{386} (2021), 1--55.

\bibitem{DSNear}J. Dong and A. Schopieray, \emph{Near-integral fusion}, arXiv:2407.15955v1 (2024).

\bibitem{DP5}J. Dong and S. Palcoux, \emph{Rank-five Grothendieck rings through multiplicity 17}, manuscript and exact-certificate supplement, 18 September 2026. The named finite-indicator and quadratic parity theorems, and Theorem 1.1, refer to this version.

\bibitem{DP4}J. Dong and S. Palcoux, \emph{Rank-four Grothendieck rings through multiplicity 32}, manuscript and exact-certificate supplement, 17 September 2026. References to numbered results are to this version.

\bibitem{EM}C. Edie-Michell, \emph{A complete classification of unitary fusion categories tensor generated by an object of dimension the golden ratio}, arXiv:1904.08909 (2019).

\bibitem{ENO}P. Etingof, D. Nikshych and V. Ostrik, \emph{On fusion categories}, Ann. of Math. (2) \textbf{162} (2005), 581--642; arXiv:math/0203060.

\bibitem{EGNO}P. Etingof, S. Gelaki, D. Nikshych and V. Ostrik, \emph{Tensor Categories}, Mathematical Surveys and Monographs \textbf{205}, American Mathematical Society, 2015.

\bibitem{ENOWGT}P. Etingof, D. Nikshych and V. Ostrik, \emph{Weakly group-theoretical and solvable fusion categories}, Adv. Math. \textbf{226} (2011), 176--205; arXiv:0809.3031v2.

\bibitem{EG}D. E. Evans and T. Gannon, \emph{The exoticness and realisability of twisted Haagerup--Izumi modular data}, Comm. Math. Phys. \textbf{307} (2011), 463--512; arXiv:1006.1326.

\bibitem{GHR}C. Galindo, S.-M. Hong and E. C. Rowell, \emph{Generalized and quasi-localizations of braid group representations}, arXiv:1105.5048v1 (2011).

\bibitem{GNNrel}S. Gelaki, D. Naidu and D. Nikshych,
\emph{Centers of graded fusion categories}, Algebra Number Theory
\textbf{3} (2009), 959--990; arXiv:0905.3117, especially Remark~2.4.

\bibitem{GNgt}S. Gelaki and D. Naidu, \emph{Some properties of group-theoretical categories}, J. Algebra \textbf{322} (2009), 2631--2641; arXiv:0709.4326.

\bibitem{GS}P. Grossman and N. Snyder, \emph{Quantum subgroups of the Haagerup fusion categories}, Comm. Math. Phys. \textbf{311} (2012), 617--643; arXiv:1102.2631v2.

\bibitem{GSAH}P. Grossman and N. Snyder, \emph{The Brauer--Picard group of the Asaeda--Haagerup fusion categories}, Trans. Amer. Math. Soc. \textbf{368} (2016), 2289--2331; arXiv:1202.4396v2, Table~1.

\bibitem{Izumi}M. Izumi, \emph{A Cuntz algebra approach to the classification of near-group categories}, arXiv:1512.04288v2 (2015), especially Example~13.8 and Section~13.3.

\bibitem{IMP}M. Izumi, S. Morrison and D. Penneys, \emph{Fusion categories between $\mathcal C\boxtimes\mathcal D$ and $\mathcal C*\mathcal D$}, arXiv:1308.5723v1 (2013).

\bibitem{IJMS}M. Izumi, V. F. R. Jones, S. Morrison and N. Snyder, \emph{Subfactors of index less than 5, part 3: quadruple points}, Comm. Math. Phys. \textbf{316} (2012), 531--554; arXiv:1109.3190.

\bibitem{KO}A. Kirillov Jr. and V. Ostrik, \emph{On a $q$-analogue of the McKay correspondence and the ADE classification of $\widehat{\mathfrak{sl}}_2$ conformal field theories}, Adv. Math. \textbf{171} (2002), 183--227; arXiv:math/0101219v3.

\bibitem{LPR}Z. Liu, S. Palcoux and Y. Ren, \emph{Classification of Grothendieck rings of complex fusion categories of multiplicity one up to rank six}, Lett. Math. Phys. \textbf{112} (2022), article 54; arXiv:2010.10264v5.

\bibitem{MY}Y. Meng and Z. Yu, \emph{Classification of some $\ZZ/2\ZZ\times\ZZ/2\ZZ$-quadratic fusion categories of rank 6}, arXiv:2608.06064v1, 6 August 2026.

\bibitem{NRW12}S.-H. Ng, E. C. Rowell and X.-G. Wen, \emph{Classification of modular data up to rank 12}, arXiv:2308.09670v2 (2025), with rank-six ancillary fusion tensors.

\bibitem{NS}S.-H. Ng and P. Schauenburg, \emph{Frobenius--Schur indicators and exponents of spherical categories}, Adv. Math. \textbf{211} (2007), 34--71; arXiv:math/0601012.

\bibitem{NSpiv}S.-H. Ng and P. Schauenburg, \emph{Higher Frobenius--Schur indicators for pivotal categories}, Contemp. Math. \textbf{441} (2007), 63--90; arXiv:math/0503167v2.

\bibitem{NRWW}S.-H. Ng, E. C. Rowell, Z. Wang and X.-G. Wen, \emph{Reconstruction of modular data from $SL_2(\mathbb Z)$ representations}, Comm. Math. Phys. \textbf{402} (2023), 2465--2545; arXiv:2203.14829.

\bibitem{NRW}S.-H. Ng, E. C. Rowell and X.-G. Wen, \emph{Recovering $R$-symbols from modular data}, arXiv:2408.02748v1 (2024).

\bibitem{Ostrik2}V. Ostrik, \emph{Fusion categories of rank 2}, Math. Res. Lett. \textbf{10} (2003), 177--183; arXiv:math/0203255.

\bibitem{Ocode}V. Ostrik, \emph{On formal codegrees of fusion categories}, Math. Res. Lett. \textbf{16} (2009), 895--901; arXiv:0810.3242.

\bibitem{Ostrik3}V. Ostrik, \emph{Pivotal fusion categories of rank 3}, Mosc. Math. J. \textbf{15} (2015), 373--396; arXiv:1309.4822.

\bibitem{CensusRepo}S. Palcoux, \emph{Fusion Ring Census}, public data and verification repository, \url{https://github.com/sebastienpalcoux/fusion-ring-census}, revision \filename{b86eb6dc6eeeafe5fc5f6906d7b074e0322e9726}, 23 September 2026. Rank-six tensors, counts and provenance are pinned to this revision.

\bibitem{SchGT}P. Schauenburg, \emph{A higher Frobenius--Schur indicator formula for group-theoretical fusion categories}, arXiv:1502.02906v1 (2015).

\bibitem{SchTwo}A. Schopieray, \emph{Categorification of integral group rings extended by one dimension}, J. London Math. Soc. \textbf{108} (2023), 1617--1641, \url{https://doi.org/10.1112/jlms.12789}. Lemma~3.1 is cited in its published form.

\bibitem{QGsurvey}A. Schopieray, \emph{Lie theory for fusion categories: a research primer}, arXiv:1810.09055 (2018).

\bibitem{SchNonPU}A. Schopieray, \emph{Non-pseudounitary fusion}, arXiv:2010.02958v1 (2020).

\bibitem{VS}G. Vercleyen and J. K. Slingerland, \emph{On low rank fusion rings}, J. Math. Phys. \textbf{64} (2023), 091703; arXiv:2205.15637v4, with ancillary multiplication tables.
\bibitem{GIcenters}P. Grossman and M. Izumi,
\emph{Drinfeld centers of fusion categories arising from generalized Haagerup subfactors},
Quantum Topol. \textbf{13} (2022), 593--668, doi:10.4171/QT/167.

\bibitem{GISAH}P. Grossman, M. Izumi and N. Snyder,
\emph{The Asaeda--Haagerup fusion categories}, arXiv:1501.07324,
especially Lemmas~4.15--4.16.

\end{thebibliography}
